\clearpage
\setcounter{page}{1}
\maketitlesupplementary

\section{Implementation Details}
\label{sec:supp_impl}

\subsection{State Intervals}
\label{sec:supp_states}

To test a state, the layer projects up to 1{,}500 surface samples of the state, placed on a 2.5\,cm grid, into each
later depth image. A pixel without depth or outside the image gives no evidence, and readings more than $60^\circ$ from
the sample normal do not count~\cite{nguyen2012kinect}. Thin parts, edges, and mislabeled pixels produce false
see-through readings. A sample is reliable if its posterior of belonging to the object is at least one half under a
two-component label model, fitted by expectation maximization at every round, and if it was never seen through while
the object was confirmed in place. The Beta density $p_A$ of \cref{eq:absence} shrinks toward pooled sensor
statistics, and an observation contributes only if it judges at least $\min(30, N_{\mathrm{rel}})$ samples.

For the second group, with the fraction $\phi$ of $N$'s surface cells that lie farther than 10\,cm from $G_A$ and their
effective number $n_{\mathrm{eff}}$, $B_{M:S} = 2\,P(\theta > \frac{1}{2})$ for
$\theta \sim \mathrm{Beta}(\frac{1}{2} + n_{\mathrm{eff}}\phi,\ \frac{1}{2} + n_{\mathrm{eff}}(1-\phi))$. The
pseudo-counts of \cref{eq:prior} come from the other objects of the same class and group, and the class mean starts
from Jeffreys' reference value $\frac{1}{2}$, or from $\frac{3}{4}$ and $\frac{1}{4}$ for classes declared movable or
fixed. Every action in a change-detection round reads $q$ from the history before that round, so the decisions are
independent of the identifiers, their order, and the end of a run.

A state $A$ ends without a successor when $\max(\operatorname{logit} q_A + L^{\mathrm{pre}}_t,\ L_t) > \ln 99$ and the
same rays see through $G_A$ more often than they support it. Here $L^{\mathrm{pre}}_t$ is the sum of \cref{eq:absence}
for a change before the state's first scored observation in the run, without the clamp at zero, and the prior term
applies only to a state inherited from the previous run. $N$ refines $A$ when $A$ has current support and
$1 - q \ge q\,B_{M:S}$. The map shows $G_A \cup G_N$ while $S$ is the most probable explanation and $G_A$ otherwise,
until a handover commits. At the end of a run, the same rules apply to the evidence gathered so far, and a remaining
$N$ that $M$ explains best goes to history. A state ended by a successor closes at the later of its last support and
the decision, and a state ended without one closes midway between its last evidence and the decision.

\subsection{Backends}
\label{sec:supp_backends}

\paragraph{TSDF.}
When a round changes the state intervals, the backend rebuilds the affected blocks by integrating their frames again in
their original order, which keeps the weight saturation of the native rule. Each face takes the identity of the nearest
current object face within one voxel and is background otherwise. Where the run's frames leave gaps in the present
surface, the backend fills them with the background faces of the active window, which it archives at every round. Fill
faces and the background faces of the previous run's shown map pass \cref{eq:continue}, the former with
$K_{\mathrm{prior}} = 0$ and the latter with $K_{\mathrm{prior}} = 1$. Background faces of the previous map leave
together with a closed object state whose old site they cover, and fill faces measured no later than that state's last
support leave as well. Object faces follow their state. A face of a continuing inherited state also leaves if it lies
inside the object's present surface, which the backend tests by casting rays in 64 fixed directions and comparing hits
on the inner and outer sides. The end of a run only drains the queues and saves the state, so the final map is the
online map.

\paragraph{Gaussians.}
FlashSplat~\cite{shen2024flashsplat} assigns identities to the Gaussians, accumulated over keyframes. The continuation
test runs at each keyframe of the current run, before GaME adds new Gaussians, with $h_e$ three standard deviations and
$\tau_e$ at least the 5\,cm tolerance. A reading covers a Gaussian only if it comes from the side from which the
previous run saw the Gaussian. A Gaussian that the run supports or covers is \emph{present} and is trained by the
current run's keyframes; one that it supports but never covers is also \emph{shared} and is trained by the keyframes of
both runs; every other Gaussian is \emph{held}: the current run leaves it unchanged, and the keyframes of its own run
keep maintaining it. Each optimization step draws one keyframe. A keyframe of the current run has weight one. An
earlier keyframe has as weight the fraction of the Gaussians it shows, among those the current run has decided on, that
are shared and that the current run has observed less than a single earlier view does: their observed area $(f/z)^2$,
summed over the views of the current run, stays below its mean over the earlier views. The earlier views of one place
thus share that place's budget, and a distant new view does not replace close earlier detail. A step updates only the
Gaussians its keyframe may train, and its loss counts only the pixels whose first Gaussian, where the transmittance
along the ray first drops below one half, is one of them. GaME's opacity pruning removes only Gaussians created in the
current run. The refinement after the last run follows the same rules and adds no Gaussians. \Cref{sec:supp_gaussian}
switches these choices off one at a time.

GaME's depth loss compares the rendered depth $D = \sum_i w_i z_i$ with the reading $d$, where $w_i$ is the blending
weight of the $i$-th Gaussian along the ray and $z_i$ is its depth. Where the accumulated opacity $A = \sum_i w_i$ is
below one, Gaussians behind the surface can satisfy this loss without contributing color, and other viewpoints then see
them in front of the surface. For a map that inherits an earlier map, we separate the two requirements of the depth
loss:
\begin{equation}
  \mathcal{L}_{\mathrm{depth}} = A\,\big|D/A - d\big| + \tau_0\,(1 - A),
  \label{eq:depth}
\end{equation}
with $A$ held fixed in the first term and $\tau_0 = 5$\,cm. The first term places the expected depth at the reading,
and the second asks the ray to end there.

\section{GaME as a Baseline}
\label{sec:supp_game}

We run GaME~\cite{yugay2026game} from its released code with seven fixes, in its own rows and under our layer alike.
(C1) The keyframe test uses the relative pose between the last keyframe and the current frame, instead of
world-to-camera poses composed in an order that depends on where a dataset places its origin. (C2) The rotation test
uses the true relative rotation, spaced by the angular radius of the calibrated image, instead of comparing Euler
angles in radians with 50. (C3) The keyframe sampler terminates when occlusion masks cover every keyframe. (C4) The
training loss counts only pixels with measured depth, whereas the released code trains the depth of a pixel without a
reading toward zero. (C5) The addition test also counts only pixels with measured depth. (C6) The refinement raises the
degree of the spherical harmonics every 1{,}000 iterations up to the declared 3, whereas the released code declares
degree 3 but never raises it. (C7) The refinement decays the position learning rate on the schedule its arguments
declare, from $1.6\cdot10^{-4}$ to $1.6\cdot10^{-6}$ over 30{,}000 steps, which the released code never applies. On Flat, the released
code reaches 24.86\,dB after refinement and GaME with these fixes 26.17\,dB.

\section{Design Choices of the Gaussian Backend}
\label{sec:supp_gaussian}

\Cref{tab:supp_gaussian} switches off one design choice of the Gaussian backend at a time. Within a pair, both runs
share the code and the first-run map. Letting the layer choose the training keyframes raises
the geometry F1 of real run B by 3.07 points and the memory tolerance by 0.64 points, and limiting the pruning keeps
681{,}011 earlier elements that GaME's pruning deletes.

\begin{table}[h]
\centering\small
\caption{One design choice of the Gaussian backend switched off at a time.}
\label{tab:supp_gaussian}
\setlength{\tabcolsep}{3pt}
\begin{tabular}{llcc}
\toprule
Design choice & Measure & With & Without \\
\midrule
Pruning limited to the current run & real B, earlier elements lost & 0 & 681{,}011 \\
                                   & real B, geometry F1 & 92.66 & 92.51 \\
                                   & simulated B, geometry F1 & 88.86 & 88.05 \\
Training pool of the layer, not every stored keyframe & real B, geometry F1 & 92.48 & 89.41 \\
Memory tolerance $\max(3\sigma, 5\,\text{cm})$, not $3\sigma$ & real B, geometry F1 & 92.99 & 92.35 \\
\bottomrule
\end{tabular}
\end{table}

\section{Refinement after the Last Run}
\label{sec:supp_refinement}

The refinement that follows the last run applies the mapping's rules: each element is trained only by keyframes of the
run that last measured it, ended elements stay fixed, and the refinement adds no elements. It raises PSNR on all three
datasets, by 3.85, 1.58, and 0.93\,dB, and changes geometry by at most 0.12 points (\cref{tab:supp_refinement}). GaME
refines one current map on every stored keyframe, including keyframes whose content has since changed, and its geometry
falls on Flat and on the real room.

\begin{table}[h]
\centering\small
\caption{Before $\to$ after the refinement that follows the last run.}
\label{tab:supp_refinement}
\setlength{\tabcolsep}{3pt}
\begin{tabular}{llcc}
\toprule
Data & Metric & GaME & Ours \\
\midrule
Flat, run 2 & PSNR & 23.42 $\to$ 26.17 & 23.79 $\to$ 27.64 \\
            & Geometry F1 & 87.71 $\to$ 84.60 & 89.78 $\to$ 89.82 \\
Simulated B & PSNR & 22.13 $\to$ 22.09 & 23.32 $\to$ 24.90 \\
            & Geometry F1 & \TODO{score} & 96.59 $\to$ 96.51 \\
Real C      & PSNR & 18.29 $\to$ 18.94 & 18.97 $\to$ 19.90 \\
            & Geometry F1 & 90.75 $\to$ 89.24 & 95.04 $\to$ 95.16 \\
\bottomrule
\end{tabular}
\end{table}

\section{Sensitivity of the State Intervals to Their Constants}
\label{sec:supp_sensitivity}

We replayed the logged evidence of the out-of-view and between-run test with each constant halved and doubled
(\cref{tab:supp_sensitivity}). The cold-start moments, the clamps, and the floors change at most two decisions. The
number of samples required in view, the sample cap, and the history gate change 4 to 12 decisions, with corrections
and new errors in both directions.

\begin{table}[h]
\centering\small
\caption{Decisions of the state test that change when one constant is halved or doubled (real B and C, simulated A
and B, replayed from the logs).}
\label{tab:supp_sensitivity}
\setlength{\tabcolsep}{3pt}
\begin{tabular}{lccc}
\toprule
Constant & Value & $\times 0.5$ & $\times 2$ \\
\midrule
Samples required in view & 30 & 9 & 12 \\
Sample cap & 1{,}500 & 12 & -- \\
History gate & 3 & 4 & 11 \\
Pseudo-observations of the population & 3 & 1 & 2 \\
Cold-start mean and variance & 0.05, 0.09 & 0 & 0 \\
Concentration floor & 2 & 0 & 2 \\
Clamps of the present fraction & 0.005, 0.995 & 0 & 0--1 \\
Clamps of the mean & 0.001, 0.999 & 0 & 0 \\
Variance floor & $10^{-4}$ & 0 & 2 \\
\bottomrule
\end{tabular}
\end{table}
\TODO{Total number of decisions in the replay.}

